\documentclass[11pt,a4paper]{article}
\usepackage[left=22mm,right=22mm,top=20mm,bottom=21mm]{geometry}
\usepackage{amsmath,fontspec,unicode-math}
\setmainfont{lmroman10-regular.otf}[
  Path=../fonts/, BoldFont=lmroman10-bold.otf,
  ItalicFont=lmroman10-italic.otf, BoldItalicFont=lmroman10-bolditalic.otf,
  Ligatures=TeX]
\setmonofont{lmmono10-regular.otf}[Path=../fonts/]
\setmathfont{latinmodern-math.otf}[Path=../fonts/]
\usepackage{booktabs,tabularx,array,graphicx,float,caption}
\usepackage{xcolor,fancyhdr}
\usepackage[numbers,sort&compress]{natbib}
\usepackage[colorlinks=true,linkcolor=navy,citecolor=navy,urlcolor=navy]{hyperref}
\definecolor{navy}{HTML}{13334F}
\definecolor{teal}{HTML}{00664F}
\definecolor{muted}{HTML}{47586A}
\definecolor{pale}{HTML}{EBF3F7}
\definecolor{ruleline}{HTML}{C9D9E2}
\hypersetup{pdftitle={Robust Sparse Smooth PCA for Handwritten Digit Recognition},pdfauthor={Paco, Sou Chak Tong}}
\newcommand{\doi}[1]{doi: \href{https://doi.org/#1}{\nolinkurl{#1}}}
\setlength{\parindent}{0pt}
\setlength{\parskip}{5pt}
\setlength{\headheight}{14pt}
\setlength{\fboxsep}{8pt}
\setlength{\emergencystretch}{2em}
\captionsetup{font=small,labelfont={bf,color=navy},justification=raggedright,singlelinecheck=false}
\graphicspath{{../../results/figures/}}
\pagestyle{fancy}
\fancyhf{}
\lhead{\small\color{muted}CISC3024 \enspace / \enspace AI Assignment \#3}
\rhead{\small\color{muted}November 2026}
\rfoot{\small\color{navy}\thepage}
\renewcommand{\headrulewidth}{0.3pt}
\newcommand{\reportfigure}[3]{%
  \begin{figure}[H]\centering
  \includegraphics[width=\linewidth]{#1}
  \caption{#2}\label{#3}
  \end{figure}}
\newcommand{\sourcequote}[3]{%
  \par\noindent{\small\texttt{#1}}\\
  \noindent\textit{\textquotedblleft#2\textquotedblright}\\
  \noindent #3\par}
\newcommand{\E}{\mathcal{E}}
\newcolumntype{Y}{>{\raggedright\arraybackslash}X}
\newcolumntype{P}[1]{>{\raggedright\arraybackslash}p{#1}}
% Generated by python -m src.summarize.
\newcommand{\maintablerows}{%
Pixels & 64 & 97.83 & 97.75\,$\pm$\,0.21 & 97.26\,$\pm$\,0.38 & 97.24 & 97.26 \\
PCA & 8 & 94.55 & 92.95\,$\pm$\,0.40 & 91.13\,$\pm$\,0.47 & 91.07 & 91.04 \\
LDA & 8 & 95.05 & 93.98\,$\pm$\,0.32 & 91.89\,$\pm$\,0.51 & 91.84 & 91.81 \\
No smoothness & 8 & 94.60 & 92.92\,$\pm$\,0.54 & 91.09\,$\pm$\,0.75 & 91.03 & 90.97 \\
RSSPCA & 8 & 94.77 & 93.12\,$\pm$\,0.58 & 90.76\,$\pm$\,0.59 & 90.70 & 90.64 \\
}
\newcommand{\validationrows}{%
0.001 & 94.72 & 94.81 & 94.39 & 90.94 \\
0.01 & 94.79 & 94.83 & 94.46 & 91.05 \\
0.1 & 93.43 & 93.81 & 94.24 & 91.65 \\
}
\newcommand{\classrecallrows}{%
PCA & 75.52 & 71.26--79.31 & 80.11 & 74.44--83.89 \\
LDA & 78.05 & 72.99--81.03 & 82.00 & 78.33--83.89 \\
No smoothness & 73.45 & 70.69--76.44 & 80.44 & 73.33--84.44 \\
RSSPCA & 72.99 & 67.82--77.59 & 79.33 & 73.89--85.56 \\
}

\begin{document}

{\bfseries\color{teal} CISC3024 --- Pattern Recognition \qquad AI Assignment \#3\par}
\vspace{2mm}
{\raggedright\bfseries\color{navy}
{\fontsize{29}{33}\selectfont Robust Sparse Smooth PCA\par}
{\fontsize{23}{28}\selectfont for Handwritten Digit Recognition\par}}
\vspace{2mm}
{\large Paco, Sou Chak Tong \enspace \textbf{U-C3-2652-3}}\par
{\small November 2026 \quad \textcolor{muted}{Initial technical report}}\par
\vspace{2mm}
{\color{ruleline}\hrule}
\vspace{3mm}
\textbf{Abstract.}
Robust sparse smooth PCA (RSSPCA) adds a pixel-grid penalty to an absolute-score criterion. This report adapts Wang et al.'s 2025 method to UCI Optdigits and evaluates it on 1,797 images from writers excluded from development. With eight coordinates and Euclidean 3-nearest neighbors, RSSPCA scored 94.77\% on clean training, 93.12\% at 25\% corrupted training and 90.76\% at 50\%. Its mean difference from PCA changed from $+0.17$ to $-0.37$ percentage points across the corrupted conditions. The grid penalty reduced loading roughness relative to the matched ablation. Fisher LDA led at eight coordinates.

\section{How AI Found the Algorithm}
The handout provided the six report topics and required AI assistance. \texttt{AGENTS.md} defined the persistent evidence and leakage rules. The separate stage file \texttt{01\_RESEARCH\_IMPLEMENT\_REPORT.md} specified the execution sequence and deliverables. Codex used the three sources to trace requirements and decisions.

\sourcequote{AGENTS.md}{Do not consult the final test set during this gate.}
{Codex downloaded only the development CSV. The independent arithmetic check and dataset smoke test ran before the selection record was written.}
\sourcequote{01\_RESEARCH\_IMPLEMENT\_REPORT.md}{do not begin the full implementation until the mandatory algorithm-selection gate (Section B2) is passed.}
{Codex screened primary papers, inspected the available reference solver and recorded a dated pass before creating the main pipeline.}
\sourcequote{AGENTS.md}{Never select algorithms, hyperparameters, figure examples or conclusions by repeatedly consulting the final test set.}
{The study fixed the example-selection rules and froze all fitted projections before obtaining the official test CSV. The second execution checked the same protocol and predictions.}

\reportfigure{workflow}{Instruction-to-execution path observed in Prompt 01. The gate links the research decision to a runnable mathematical and dataset probe.}{fig:workflow}

\clearpage
\subsection{Search, screening and selection}
Codex ran a targeted search on 17 November 2026. The discovery criteria were a 2024--2026 original peer-reviewed method, a substantive PCA or LDA contribution, direct vision relevance and enough detail for a controlled Python evaluation. The initial web searches and solver follow-up used these exact strings.

{\small\setlength{\tabcolsep}{5pt}
\begin{tabularx}{\linewidth}{@{}P{21mm}Y@{}}
\toprule
Purpose & Query string \\
\midrule
RSSPCA & \texttt{2025 robust sparse smooth principal component analysis RSSPCA face recognition} \\
Sparse LDA & \texttt{2024 robust and sparse linear discriminant analysis face image objective robust sparse LDA} \\
2D methods & \texttt{2025 2026 robust two dimensional principal component analysis face recognition algorithm sparse Lp} \\
Broader search & \texttt{2024 2025 novel two dimensional PCA image recognition generalized weighted 2DPCA face recognition open access algorithm} \\
LDA full text & \texttt{"Towards robust and sparse linear discriminant analysis for image classification" pdf RSLDA+ Liu Feng} \\
Code & \texttt{DLR-2DPCA Pengfei Bi GitHub double layer robust code} \\
\bottomrule
\end{tabularx}}\par
A literature lookup recovered the SCRAMBLE preprint. Codex used the journal article to verify the method and wear-scar image application. The complete query ledger is in \path{docs/research.md}. The literature-review skill guided primary-source screening. The computation-audit skill guided the independent calculation. The experimental-design and statistical-analysis skills informed the shared corruption masks and the interpretation of variation across runs \citep{kassis2026}.

\begin{table}[H]\small
\caption{Shortlist and internal suitability scores. Weights for recency, PCA/LDA centrality, implementation, experimental fairness, mathematical verification and figure potential were 15/20/20/20/15/10 percent. These are project decision weights.}
\begin{tabularx}{\linewidth}{@{}P{26mm}YP{56mm}r@{}}
\toprule
Candidate & Contribution and vision use & Decision evidence & /5\\
\midrule
RSSPCA \citep{wang2025} & Absolute projected dispersion, sparse and grid-smooth loadings. Face recognition and reconstruction. & Open equations and MATLAB reference. Direct smoothing ablation at fixed feature count. & 5.00\\
DLR-2DPCA \citep{bi2025} & Cut reconstruction loss and adaptive weights. Face, digit and traffic sign recognition. & Readable algorithm. More numerical boundary cases. Matrix outputs restrict matched budgets. & 4.30\\
SCRAMBLE \citep{pfeiffer2025} & Robust cellwise reconstruction and elastic net. Wear-scar image regression. & Open R reference. More solver and scale choices for transfer to digit recognition. & 4.15\\
RSLDA+ \citep{liu2024} & Sparse LDA with ADMM and thresholding. Image classification. & Complete solver unavailable through sources accessed during the search. & 2.95\\
\bottomrule
\end{tabularx}
\end{table}
\textbf{Decision gate.}
Wang et al. published RSSPCA in PLOS One in May 2025. Its four-neighbor graph gives a concrete mechanism to test on raster data. At a fixed coordinate budget, the graph penalty can be switched off while keeping the absolute-score objective, sparsity setting and classifier. This made it possible to compare the effect of spatial regularization separately from the overall comparison with PCA and LDA.

\clearpage
\subsection{Why RSSPCA passed the gate}

Recency and a direct change to PCA or LDA were necessary conditions. For the remaining candidates, the question was whether the published equations could support an independently checked implementation and a fair experiment. Table~\ref{tab:gate-scores} shows the individual ratings behind the shortlist.

\begin{table}[H]\small\centering
\caption{Recorded gate ratings from 1 to 5. A weighted score averages the six ratings with the indicated weights. Eligibility also requires an accessible solver.}
\label{tab:gate-scores}
\setlength{\tabcolsep}{5pt}
\begin{tabular}{@{}lrrrrrrr@{}}
\toprule
Candidate & Recency & Centrality & Feasibility & Fairness & Math & Figures & Score\\
Weight (\%) & 15 & 20 & 20 & 20 & 15 & 10 & \\
\midrule
RSSPCA & 5 & 5 & 5 & 5 & 5 & 5 & 5.00\\
DLR-2DPCA & 5 & 5 & 3 & 4 & 5 & 4 & 4.30\\
SCRAMBLE & 5 & 5 & 3 & 4 & 4 & 4 & 4.15\\
RSLDA+ & 5 & 5 & 1 & 2 & 2 & 3 & 2.95\\
\bottomrule
\end{tabular}
\end{table}

DLR-2DPCA was a viable alternative with a readable algorithm and direct digit-recognition precedent. On the $8\times8$ grid, a right projection with $r$ columns would deliver $8r$ scalar scores, which restricts the dimensions that can be matched to vector PCA and LDA. SCRAMBLE also had an accessible method and reference package. Its wear-scar application used regression, and transferring the method here would require choices about the cellwise loss, scale and sparsity fitting. RSLDA+ addressed image classification directly, but the accessible material lacked the complete iteration. The scores reflect those practical differences and carry no information about test accuracy.

\subsection{What the development-only probe measured}

The two-pixel calculation independently solved the defining linear system by Cramer's rule. The numerical residual was $6.28\times10^{-16}$ and the smallest eigenvalue of the symmetric system was $0.873$. A separate calculation with an exactly zero loading coordinate checked the branch most exposed to the inverse-diagonal formulation. The normalized update and its stopping interpretation are derived in Section~\ref{sec:method}.

The data check read only the 3,823-row development file. It found 3,823 unique image vectors, ten classes with 376 to 389 examples each, and two constant coordinates. On the 3,058/765 feasibility split, PCA with eight coordinates and Euclidean 3-NN reached 96.08\%. A one-direction RSSPCA probe met the $10^{-4}$ step threshold in 27 iterations. These checks established that the representation, evaluator and update could run on the available CPU before the comparison study was designed. The 96.08\% value refers to that feasibility split.

The dated record in \path{docs/algorithm_selection.md} locks the method and adaptations before the main implementation. The later study lock in \path{results/selection.json} records the dimension and selected penalties. Keeping the two records distinct lets a reviewer trace the decision to use RSSPCA and the subsequent selection of its parameters.

\clearpage
\section{Algorithm Description}
\label{sec:method}
\subsection{Classical reference subspaces}

Let $a_i\in\mathbb R^d$ be training image $i$, flattened in row-major order, with digit $y_i\in\{0,\dots,C-1\}$. There are $n$ training images, $d=64$ coordinates and $C=10$ classes. The training mean is $\mu=n^{-1}\sum_i a_i$. Let $X\in\mathbb R^{n\times d}$ have rows $x_i^\top=(a_i-\mu)^\top$.

Classical PCA \citep{pearson1901} solves $\max_{\|w\|_2=1} w^\top S w$ for $S=X^\top X/(n-1)$. Differentiating with the norm constraint gives $Sw=\lambda w$. With thin SVD $X=U\Sigma V^\top$, the first $k$ eigenvectors are the first $k$ columns of $V$, with eigenvalues $\sigma_j^2/(n-1)$. The rank is at most $\min(d,n-1)$, and the experiment uses unwhitened scores.

For Fisher LDA \citep{fisher1936}, let $\mu_c$ and $n_c$ be the class mean and size. The within- and between-class scatters are
\[
 S_W=\sum_{i=1}^{n}(a_i-\mu_{y_i})(a_i-\mu_{y_i})^\top,\qquad
 S_B=\sum_{c=0}^{C-1}n_c(\mu_c-\mu)(\mu_c-\mu)^\top.
\]
When $S_W$ is positive definite, maximizing $w^\top S_B w$ under $w^\top S_Ww=1$ yields $S_Bw=\lambda S_Ww$. The weighted class means sum to zero after global centering, so $\operatorname{rank}(S_B)\le C-1=9$. The rank-aware SVD fit works in the nonzero within-class directions. This handles constant pixels without forming $S_W^{-1}$. Digit labels inform the LDA subspace. The downstream classifier is shared across all subspaces.

\subsection{RSSPCA objective and the image graph}

For extraction step $j$, let $R\in\mathbb R^{n\times d}$ be the residual training matrix, with $R=X$ initially. Wang et al.\ \citep{wang2025} define the loading objective
\begin{equation}
 \max_w \|Rw\|_1
 \quad\text{subject to}\quad
 \|w\|_2=1,\qquad \|w\|_1\le c_1,\qquad w^\top Lw\le c_2 ,
 \label{eq:objective}
\end{equation}
where $c_1,c_2$ are sparsity and roughness bounds. Absolute projected scores grow linearly with projection magnitude, compared with quadratic growth in PCA. This follows the earlier PCA-$\ell_1$ criterion \citep{kwak2008}. For example, doubling one residual row at fixed $w$ doubles its absolute-score contribution and quadruples its squared-score contribution. A large residual may still affect the direction because the sign vote multiplies the residual row.

Build $A$ by joining horizontal and vertical neighbors and use $L=\operatorname{diag}(A\mathbf1)-A$. In an $8\times8$ grid, there are 112 undirected edges and no wrap around at a row boundary. For edge set $\E$,
\begin{equation}
 w^\top Lw=\tfrac{1}{2}\sum_{p,q}A_{pq}(w_p-w_q)^2
           =\sum_{\{p,q\}\in\E}(w_p-w_q)^2 .
 \label{eq:energy}
\end{equation}
Thus $L$ is symmetric positive semidefinite and annihilates constant vectors. Penalizing this energy favors nearby pixel weights that agree. The $\ell_1$ penalty on unit vectors favors concentration on fewer pixels. These effects can compete near sharp stroke boundaries.

\reportfigure{solver}{One extraction step and the shared nearest-neighbor evaluator. The grid Laplacian acts on loading coordinates. Deflation uses training residuals, whereas digit predictions use scores from the centered, original input.}{fig:solver}

\clearpage
\subsection{From the minorizer to the published update}

For a current unit direction $w$ with nonzero coordinates, set $D=\operatorname{diag}(|w|)$ and $v=R^\top\operatorname{sign}(Rw)\in\mathbb R^d$, using $\operatorname{sign}(0)=0$. For any trial $z\in\mathbb R^d$,
\[
 \|Rz\|_1\ge v^\top z,\qquad
 |z_p|\le\frac{z_p^2}{2|w_p|}+\frac{|w_p|}{2}.
\]
Both inequalities are tight at $z=w$. The first follows termwise from $|b|\ge s b$ for $|s|\le1$. The second follows by expanding $(|z_p|-|w_p|)^2\ge0$.

Write the Lagrangian with penalty multipliers $a,b\ge0$ on the last two constraints in Eq.~\eqref{eq:objective} and $c>0$ on $\|z\|_2^2-1$. Replacing the two absolute-value terms with their bounds gives the concave quadratic minorizer, up to an additive constant,
\[
 G(z\mid w) = v^\top z -\tfrac a2 z^\top D^{-1}z -b z^\top Lz-c z^\top z.
\]
Its unconstrained stationary point solves $(2cI+aD^{-1}+2bL)z=v$. The published update normalizes the direction, absorbing the positive scale $2c$ by defining $\eta=a/(2c)$ and $\gamma=b/c$. The pre-normalization vector $p$ solves
\begin{equation}
 (I+\eta D^{-1}+\gamma L)p=v,\qquad
 w^+=p/\|p\|_2 .
 \label{eq:original}
\end{equation}
For invertible $D$, putting $p=Du$ gives the published zero-safe form $p=D(D+\eta I+\gamma LD)^{-1}v$. Our implementation uses $H=D^{1/2}$ and the equivalent symmetric system
\begin{equation}
 M t=Hv,\quad M=D+\eta I+\gamma H L H,\qquad
 p=Ht .
 \label{eq:spd}
\end{equation}
The identity $H(D+\eta I+\gamma LD)=MH$ ensures equivalent outputs even with zero entries in $D$ when $\eta>0$. The nonsymmetric system is invertible because it is similar to the positive-definite block on the nonzero coordinates and has diagonal $\eta I$ in its remaining block. Also $M\succeq\eta I$. The grid has maximum degree four, so $\|L\|_2\le8$. With unit $w$, \(\kappa_2(M)\le(\eta+1+8\gamma)/\eta\). For the selected penalties this bound is 125. The implementation solves the system using a Cholesky-based routine.

\subsection{Worked example and stopping interpretation}

Take rows $(1,2)$ and $(-1,-2)$, $w=(0.6,0.8)^\top$, and a two-pixel edge with $L=\left[\begin{smallmatrix}1&-1\\-1&1\end{smallmatrix}\right]$. Then $v=(2,4)^\top$. For $\eta=0.2$ and $\gamma=0.5$, the nonsymmetric system for $u$ has rows $(1.1,-0.4)$ and $(-0.3,1.4)$. Cramer's rule followed by $p=Du$ gives
\[
 p=(132/71,\,200/71)^\top,\qquad w^+=(0.550842,\,0.834609)^\top.
\]
The SPD computation agreed. Starting with $(1,0)^\top$ gave $p=(20/17,0)^\top$, illustrating that exact zero coordinates are absorbing for $\eta>0$.

The method maximizes the unconstrained quadratic before normalization. A sphere-constrained quadratic optimizer would also satisfy tangent stationarity. The numerical example gives $\|(I-w^+{w^+}^\top)(v-Qw^+)\|_2=0.376$ for $Q=I+\eta D^{-1}+\gamma L$, so the normalized candidate fails that stationarity condition. We use the published map as a fixed-point iteration. A small direction change serves as its operational stop. Exact constraint-budget satisfaction, monotonic ascent of the absolute-score objective and global optimality remain unverified.

\clearpage
\section{How AI Implemented the Algorithm}

Codex inspected the authors' GPL-3.0 MATLAB reference \citep{matlab} and implemented the paper equations independently with NumPy and SciPy. The reference initializes with axes from the original SVD, stops using a relative change in the L2 score norm and uses accumulated raw directions in deflation. This project initializes each component with the leading covariance eigenvector of the \emph{current residual}. It uses the sign-invariant step size $\min(\|w^+-w\|_2,\|w^++w\|_2)$ and stops at $10^{-4}$ or 100 updates. The code records both step size and mean absolute score. The largest-magnitude entry is made nonnegative for display. With equal leading eigenvalues, the initialization can depend on the numerical basis picked by the solver.

\subsection{Deflation, coordinate scores and cost}

Write $W_j=[w_1,\ldots,w_j]$ and let $Q_j$ have orthonormal columns spanning $W_j$, maintained by twice removing the existing basis projection from each new column before normalizing. We recompute \(R_j=X-(XQ_j)Q_j^\top\). Then $R_j W_j=0$ up to roundoff. The expression $I-W_jW_j^\top$ is an orthogonal projector only when the columns of $W_j$ are orthonormal. The raw columns can remain oblique, so prediction uses $W_k^\top(a-\mu)$ and retains the learned metric on those coordinates. The implementation rejects excess dimensions, exhausted residuals, dependent new directions and an update with zero norm.

For two images $a$ and $b$, squared Euclidean distance between the raw scores is
\[
 \|W_k^\top(a-b)\|_2^2=(a-b)^\top W_kW_k^\top(a-b).
\]
The common mean cancels in this calculation. Flipping the sign of a loading leaves this distance unchanged. Using the orthonormal basis $Q_k$ for classification would instead assign equal weight to all directions in the learned span and could change which neighbors vote. The two bases therefore have distinct roles in the implementation.

For $k$ components and $T$ updates per component, a dense implementation at $n\ge d$ costs $O\!\left(k(nd^2+d^3)+kT(nd+d^3)+ndk^2\right)$ time. The terms bound residual covariance/eigenvector initialization, sign votes and SPD solves, and growing orthogonal projections. Working storage is $O(nd+d^2+dk)$ plus diagnostic traces. A test score costs $O(dk)$ per image and brute-force 3-NN compares $k$-coordinate scores against $n$ training examples. The $64\times64$ system is small enough for dense solves on this dataset.

\begin{table}[H]\small
\caption{Pseudocode and equation-to-code mapping. Fit inputs are training rows, labels, penalties and the feature budget. Each returned vector has 64 entries.}
\begin{tabularx}{\linewidth}{@{}P{9mm}YP{53mm}@{}}
\toprule
Step & Operation & Implementation \\
\midrule
1 & Compute $\mu$, center training rows, set $R\gets X$, and form $L$. & \texttt{fit\_rsspca}, \texttt{grid\_laplacian} \\
2 & For each component, initialize $w$ from the leading residual covariance eigenvector. & \texttt{fit\_rsspca} \\
3 & Form $v,D,H,M$, solve Eq.~\eqref{eq:spd}, normalize, and repeat to the step threshold or cap. & \texttt{fixed\_point}, \texttt{iterate\_component} \\
4 & Append $w$ to $W$ and reproject $X$ away from its span. & \texttt{fit\_rsspca} \\
5 & Project train and test with the saved $\mu,W_k$. Vote among the 3 closest training scores. & \texttt{Projection.transform}, \texttt{classifier} \\
\bottomrule
\end{tabularx}
\end{table}

\clearpage
\subsection{Verification and execution}

The verification suite tests mathematical invariants against independent references and exercises the input boundaries that could silently change an experiment. All ten collected tests passed. The small development smoke run used 300 fitting rows and 50 validation rows, exercising all five methods before the complete run.

\begin{table}[H]\small
\caption{Representative verification checks from \texttt{tests/test\_pipeline.py}. Numerical tolerances are absolute floors used by the assertions.}
\label{tab:checks}
\begin{tabularx}{\linewidth}{@{}P{33mm}Y@{}}
\toprule
Check & Evidence and the failure it can expose\\
\midrule
Graph geometry & A $2\times3$ grid has seven unique edges, zero row sums, a nonnegative spectrum up to roundoff and no edge spanning a row boundary. A hand-summed edge energy checks the flattening convention.\\
Update equivalence & The SPD solve is compared with an independent nonsymmetric solve at three penalty pairs, including $\eta=10^{-9}$ with $\gamma=0$ and $\eta=10^{-8}$ with $\gamma=100$. The test uses a $10^{-10}$ absolute tolerance floor after aligning signs and verifies that exact zero coordinates remain zero.\\
Deflation and prefix & Tests compare repeated loading fits and the first loading extracted on its own. An independently computed basis gives a residual orthogonal to the learned span within $10^{-13}$. A rank-deficient input must raise.\\
Classical baselines & PCA is checked against the scikit-learn subspace and the covariance eigen-equation. The Fisher SVD subspace is checked against a generalized eigenproblem and adding a constant pixel must preserve its scores up to sign.\\
Data and selection & The two corruption severities share row IDs, patch locations and values on their intersection. A serialization round trip checks mask replay. A tie in mean accuracy exercises the stated selection order.\\
\bottomrule
\end{tabularx}
\end{table}

\textbf{The boundary between fitting and test scoring.}
The validation stage checks the development-file digest, writes the folds and records the chosen parameters. The preparation stage refits using the full development file and saves the centers, loading matrices and training masks. A manifest hashes the selection, split, fitted and fit-record files. The final stage verifies those hashes and the fitting-code hashes before loading the official test file. It reconstructs the training exemplars from the saved masks and projects both partitions with the frozen centers and weights.

\textbf{Replay and metric checks.}
The executed run used Windows 11, Python 3.14.3, NumPy 2.4.4, SciPy 1.18.0 and scikit-learn 1.9.0 on CPU, with BLAS and OpenMP limited to one thread. An independent NumPy check reconstructed all confusion matrices and metrics using the saved predictions and official labels. A separate run repeated tuning, fitting and evaluation and matched all 165 prediction arrays exactly within that environment. Fit times are wall-clock measurements and can vary across runs.

\clearpage
\section{Experiment Settings and Results}
\subsection{Data, corruption and comparators}

Optdigits \citep{uci1998} counts on-pixels in $4\times4$ cells of normalized $32\times32$ digit images, giving a row-major $8\times8$ count grid. Counts were divided by 16 and the training mean was fitted anew under each corruption condition. All 64 coordinates were kept to preserve the grid. The official development file has 3,823 images from 30 writers. The official test file has 1,797 images from 13 other writers. All image vectors were unique within each file and none exactly matched across files. The source is licensed CC BY 4.0.

\reportfigure{dataset}{Development data and one synthetic training patch. A, first official development row of each class, rendered from block counts. B, counts by digit in the development file. C, the first corrupted row in the prespecified seed-7301 mask before and after replacement. The outline marks the affected $4\times4$ cells.}{fig:dataset}

At fractions 0.25 and 0.50, the study sampled $\lfloor fn\rfloor$ training rows without replacement. One $4\times4$ block per selected image was drawn uniformly from 25 valid positions. Each entry in that block was independently replaced by black or white at equal probability. The treatment targets one quarter of the selected image's grid, or about 6.25\% and 12.5\% of all training cells at the two fractions. A cell can retain its original value if it already matches the replacement. A seed determines both row order and patch assignments. The 25\% training mask is a prefix of the 50\% mask for that seed, and every method sees the same masks and labels.

Final refits used 955 or 1,911 corrupted training rows and seeds 7301--7305. Clean training ran once. Validation and test images stayed clean. The treatment affects both the estimated projection and the training exemplars available to the nearest-neighbor classifier. The raw-pixel comparator has no fitted projection criterion, so any change in its accuracy under this protocol comes from the modified exemplars.

Three shuffled stratified development folds with seed 193 were used for tuning. Fold corruption seeds were 6101--6103. Each fold fitted on 2,548 or 2,549 images and validated on the remaining 1,275 or 1,274 images. The data do not include per-row writer IDs, so the internal folds may share writers. The official test split retains the source's writer boundary.

\textbf{Comparison budget.} The fixed Euclidean 3-NN evaluator used uniform voting and the same training row order. All reduced methods used four and eight coordinates. PCA and the robust arms also used 16 and 24. Eight was the primary budget because it compresses the 64-input representation eightfold and remains below LDA's nine-direction rank ceiling. The 64-coordinate raw-pixel reference keeps the original feature budget. PCA, RSSPCA and the no-smoothing ablation fit directions without labels. LDA uses the training labels to choose a subspace, and all methods use those labels in the final voting stage.

\clearpage
\subsection{Development selection and metric definitions}

The tuning grid combined $\eta\in\{0.001,0.01,0.1\}$ with $\gamma\in\{0.03,0.3,3\}$. The no-smoothing check repeated the same three sparsity coefficients with $\gamma=0$. For each pair, the objective was mean eight-coordinate validation accuracy across nine equally weighted fold-by-corruption conditions. Equal accuracy favors smaller $\gamma$, then smaller $\eta$. Each fold fit its own training center, projection and neighbor model before evaluating its clean validation rows.

\begin{table}[H]\small\centering
\caption{Mean development accuracy (\%) over the nine fold-by-corruption conditions. Each cell summarizes three development folds and three training corruption levels. The no-smoothing check occupies the $\gamma=0$ column.}
\label{tab:tuning}
\begin{tabular}{@{}rrrrr@{}}
\toprule
$\eta$ & $\gamma=0$ & $\gamma=0.03$ & $\gamma=0.3$ & $\gamma=3$\\
\midrule
\validationrows
\bottomrule
\end{tabular}
\end{table}

Among positive graph penalties, $(0.01,0.03)$ achieved the highest mean, 94.83\%. The corresponding $(0.001,0.03)$ result was 94.81\%, a difference of only 0.02 percentage points after rounding. Every $\gamma=3$ entry was below every $\gamma=0.03$ entry. The grid therefore shows clearer sensitivity to the largest smoothing penalty than to the choice between the two leading positive-penalty settings. The no-smoothing check independently selected $\eta=0.01$ at 94.79\%, the same $\eta$ used in the matched ablation. The complete grid and fitted projections were saved before test scoring.

Setting $\gamma=0$ makes Eq.~\eqref{eq:spd} diagonal and gives $p_i=|w_i|v_i/(|w_i|+\eta)$. At a fixed sign vote, this factor attenuates coordinates with smaller current magnitudes. Positive $\gamma$ adds the off-diagonal element $-\gamma\sqrt{|w_iw_j|}$ for adjacent pixels. The sign votes and sparsity weighting are shared across both variants, which isolates the added graph term.

Let $C_{ij}$ count official test images with true digit $i$ and predicted digit $j$. For class $c$, let $n_c=\sum_j C_{cj}$ and $\hat n_c=\sum_i C_{ic}$. The reported scores are
\[
\mathrm{Accuracy}=\frac{\sum_c C_{cc}}{\sum_c n_c},\qquad
\mathrm{Balanced\ accuracy}=\frac{1}{10}\sum_c\frac{C_{cc}}{n_c},\qquad
\mathrm{Macro\!-\!F1}=\frac{1}{10}\sum_c\frac{2C_{cc}}{n_c+\hat n_c}.
\]
Overall accuracy weights images equally. Balanced accuracy gives each digit the same weight by averaging its recall. Macro-F1 also penalizes assigning too many other digits to a class. All ten digits are represented in the official split, with 174 to 183 images per class.

For the corrupted conditions, each accuracy entry averages five evaluations on the same official test images. The sample standard deviation describes variation induced by training masks. A paired contrast first subtracts the two methods' accuracies for each shared mask and then averages those differences. This comparison uses the same corrupted images in both fits and the same test examples in both evaluations. Neither the sample standard deviation nor the paired range measures sampling uncertainty across new writers.

\clearpage
\subsection{Recognition and the smoothing ablation}

\begin{table}[H]
\small
\caption{Official test performance (\%). Clean training is a single fit. Corrupted accuracy entries are mean $\pm$ sample SD across five matched training masks. Balanced accuracy and macro-F1 columns give means at 50\% corruption.}
\centering
\begin{tabular}{@{}lrrrrrr@{}}
\toprule
 & & \multicolumn{3}{c}{Accuracy by training corruption} &
 \multicolumn{2}{c}{At 50\%} \\
\cmidrule(lr){3-5}\cmidrule(l){6-7}
Method & Coord. & Clean & 25\% & 50\% & Bal.\ acc. & Macro-F1 \\
\midrule
\maintablerows
\bottomrule
\end{tabular}
\label{tab:results}
\end{table}

\reportfigure{comparison}{Held-out recognition and matched contrasts. A, method means with one point per corruption draw at 25\% and 50\%. The raw-pixel reference uses 64 coordinates. B, run-level differences in percentage points (pp) using the same mask within a pair. Diamonds mark the means.}{fig:comparison}

\textbf{RSSPCA versus classical PCA.}
The mean differences were $+0.22$ pp on clean training, $+0.17$ pp at 25\% corruption and $-0.37$ pp at 50\%. The paired ranges were $[-0.39,+1.00]$ pp and $[-1.00,+0.22]$ pp at the two corruption levels. RSSPCA led PCA on three of five masks at 25\% and one of five at 50\%. From clean to 50\% corruption, RSSPCA fell by 4.01 pp and PCA by 3.42 pp, so the absolute-score and graph penalties did not yield a smaller loss under this training perturbation.

\textbf{Smoothness at fixed sparsity.}
Adding the graph term increased accuracy by $0.20$ pp on average at 25\% corruption and decreased it by $0.32$ pp at 50\%. The paired ranges were $[-0.11,+0.67]$ and $[-1.61,+0.11]$ pp. Changing only $\gamma$ isolates the grid regularizer in the update. It also changes fitted directions and their deflation trajectory, so the comparison measures the full representation. The accuracy benefit depended on corruption severity.

\textbf{Matched supervised control.}
Fisher LDA led the eight-score comparison under all three conditions and beat RSSPCA on all five masks at 50\% corruption, by $1.12$ pp on average. That advantage is consistent with using labels to retain class-separating directions. The raw-pixel reference led overall with 64 coordinates and lost only 0.57 pp from clean to 50\% training corruption. The comparison with raw pixels therefore also reflects the cost of compressing this small image space.

\clearpage
\subsection{Feature budget, cost and loading diagnostics}

\reportfigure{sensitivity}{Feature budget and CPU fit time at 50\% corrupted training. A, mean $\pm$ sample SD on five paired corruption draws. The horizontal dotted line gives the raw-pixel mean at 64 coordinates. B, each point is a fit-time measurement at eight coordinates (64 for pixels), with medians shown by outlined markers.}{fig:sensitivity}

At 50\% corruption, RSSPCA rose from 73.72\% at four coordinates to 90.76\% at eight, 95.78\% at 16 and 97.00\% at 24. Classical PCA reached 96.92\% at 24. The 64-coordinate raw reference reached 97.26\%. The gain from retaining more coordinates was substantially larger than the difference between robust and classical PCA at the same budget. These dimensions were specified before testing, and the primary four-way reduced comparison remains at eight.

Median fit times at eight coordinates were 69.5 ms (RSSPCA), 75.1 ms (no smoothness), 10.4 ms (PCA) and 14.9 ms (LDA). Raw pixels took 2.3 ms. Fit times cover projection and 3-NN training after mask construction. Batched prediction averaged 15--16 ms at eight coordinates and 27 ms for raw pixels.

\reportfigure{loadings}{Fixed seed 7301 with 50\% corrupted training. A, the first three loadings share one signed color scale. Row labels include mean grid energy across the first eight unit directions. B, step norm by iteration for the first eight loadings of each robust arm. The dotted line marks the stopping tolerance.}{fig:loadings}

At 50\% corruption, smoothing lowered mean eight-direction grid energy from 2.041 to 1.906 (6.6\%). The share of entries with $|w_p|\le10^{-3}$ fell from 23.1\% to 20.0\%, consistent with weight spreading to neighbors. Both arms met tolerance for all 88 rank-eight directions. In the 24-coordinate fits, 2 of 264 RSSPCA directions and 5 of 264 no-smoothing directions reached the cap and were retained.

\clearpage
\subsection{Class errors and interpretation}

\reportfigure{class_errors}{Per-class behavior for the prespecified seed 7301 at 50\% corrupted training and eight coordinates. A, RSSPCA confusion matrix, with raw counts over row-normalized color. B, recall by true digit for the classical and RSSPCA projections.}{fig:errors}

RSSPCA recovered 118 of 174 eights (67.8\%) and 133 of 180 nines (73.9\%) in this draw. The largest off-diagonal counts were eight predicted as one (24) and nine predicted as three (20). LDA recalled 78.2\% of eights and 78.3\% of nines in the same draw. Table~\ref{tab:recall} puts that example in context. Mean RSSPCA recall remained lower than PCA and LDA for both digits across the five masks, although individual ranges overlap. Eight-coordinate embeddings may discard stroke details used to separate some digits. The confusion counts cannot identify which details were lost.

\begin{table}[H]\small\centering
\caption{Descriptive follow-up on recall (\%) for digits 8 and 9. Each mean and range uses five 50\% training corruption masks, the same test images and eight coordinates. The two classes are examined because they account for the largest recall deficits in Figure~\ref{fig:errors}.}
\label{tab:recall}
\begin{tabular}{@{}lrrrr@{}}
\toprule
 & \multicolumn{2}{c}{Digit 8, $n=174$} & \multicolumn{2}{c}{Digit 9, $n=180$}\\
\cmidrule(lr){2-3}\cmidrule(l){4-5}
Method & Mean & Range & Mean & Range\\
\midrule
\classrecallrows
\bottomrule
\end{tabular}
\end{table}

\textbf{Limitations.}
The study uses one low-resolution digit dataset and one official split. Validation can share writers across folds, because the public file lacks per-row writer IDs. The source-defined holdout is writer-disjoint, but the design cannot estimate how the chosen parameters would vary across other groups of writers. Synthetic patches affect both the fitted subspace and stored exemplars. A factorial training test with independently perturbed projection-fit and neighbor-store images would distinguish those mechanisms. Exact zeros, the residual start and the iteration limit can affect the learned directions. Generalization to real occlusion, higher-resolution images and other classifiers remains open.

\clearpage
\subsection*{Conclusion}

RSSPCA's grid penalty reduced loading roughness on Optdigits. Its recognition accuracy stayed close to PCA and the relative advantage changed across corruption rates. LDA led the eight-coordinate comparison. Under heavy training corruption, increasing the coordinate budget helped each PCA-family method more than changing the projection objective.

\section{What I Have Learnt from This AI Assignment}

\noindent\colorbox{pale}{\parbox{\dimexpr\linewidth-2\fboxsep\relax}{%
\textbf{Student-authored reflection pending.} Paco, Sou Chak Tong will supply the personal learning reflection before submission. The items below identify evidence in the draft that the student can draw on.}}\par

A reflection can discuss how the source and numerical checks changed trust in an AI-generated implementation, using the two-pixel update or the pre-test lock as an example. It can also connect a personal observation to the study's central result. The smoothness term behaved as intended on the loadings, while recognition still depended on the coordinate budget and training corruption. Any account of what the student learned or would do differently belongs in the student's own words.

\section{Webpage Link of the Source Codes}

\noindent\colorbox{pale}{\parbox{\dimexpr\linewidth-2\fboxsep\relax}{%
\textbf{Public GitHub URL pending.} This review draft accompanies the local source. The public link will be inserted after the repository has been published and checked.}}\par

\begin{table}[H]\small
\caption{Evidence available in the review workspace. Paths are relative to the project root.}
\begin{tabularx}{\linewidth}{@{}P{54mm}Y@{}}
\toprule
Path & Reviewer use \\
\midrule
\path{README.md}, \path{configs/study.json} &
Environment and run commands, fold and corruption seeds, feature budgets and stopping settings.\\
\path{docs/algorithm_selection.md}, \path{docs/ai_usage.md} &
Selection gate, prompt and search trail, and recorded chronology from method choice to test scoring.\\
\path{results/validation.json}, \path{results/selection.json}, \path{results/frozen.json} &
Full development scores, selected parameters, fitting-source hashes and frozen artifact digests.\\
\path{results/predictions.npz}, \path{results/final.json} &
Official test row IDs, ground truth and keyed predictions, with per-run confusion counts and timing.\\
\path{report/report.tex}, \path{src/figures.py}, \path{src/summarize.py} &
Report entry point and the commands that regenerate its figures and numeric tables from the archives.\\
\bottomrule
\end{tabularx}
\end{table}

To replay from the project root, use the \texttt{src.experiment} module to fetch the development data, then run \texttt{smoke}, \texttt{tune} and \texttt{prepare}. Once the freeze manifest exists, run \texttt{fetch-tes}, \texttt{final} and \texttt{verify} in a fresh output directory. For a report-only check, \texttt{src.summarize} reconstructs metrics from the archived predictions and official labels. The figures can be exported separately using \texttt{src.figures}. The UCI data source is linked in the references.

\clearpage
\renewcommand{\bibsection}{\section*{References}}
\bibliographystyle{plainnat}
\bibliography{references}
\end{document}
